\documentclass{HELBreport}
\usepackage{fancyhdr} %used with \pagestyle.
\usepackage[style=apa, backend=biber, citestyle=numeric]{biblatex}
%Bibliography
\DeclareLanguageMapping{french}{french-apa}
\addbibresource{biblio.bib}
\DeclareFieldFormat{labelnumberwidth}{\mkbibbrackets{#1}}

\defbibenvironment{bibliography}
  {\list
     {\printtext[labelnumberwidth]{%
      \printfield{labelprefix}%
      \printfield{labelnumber}}}
     {\setlength{\labelwidth}{\labelnumberwidth}%
      \setlength{\leftmargin}{\labelwidth}%
      \setlength{\labelsep}{\biblabelsep}%
      \addtolength{\leftmargin}{\labelsep}%
      \setlength{\itemsep}{\bibitemsep}%
      \setlength{\parsep}{\bibparsep}}%
      \renewcommand*{\makelabel}[1]{\hss##1}}
  {\endlist}
  {\item}

\sceau{Pictures/sceauHELB.png}

\frenchbsetup{StandardLists=true}

\begin{document}


\titleHELB{
	title={Titre de travail de fin d’étude},
	studies={},
	course={Travail de fin d’études présenté en vue de l’obtention du diplôme de Bachelier(ère) en Informatique de Gestion.},
	author={Nom Prénom\\ (N° de matricule\\ - voir carte étudiant)},
	date={\textcolor{gray}{© \the\year{} John Smith (Numéro de matricule)}},
	teacher={Promoteur : M. ou Mme. Nom Prénom},
	logo={Pictures/HELB-head.png},
	manyAuthor
}

\pagestyle{fancy} % page numbering at bottom of the page.
\chapter*{Remerciements}

\textit{Vos remerciements ici ...}

\tableofcontents

\chapter{Introduction}


L’introduction est en général organisée comme suit :
\begin{itemize}
    \item Présentation du contexte ;
    \item Définition du problème ;
    \item Présentation de la solution proposée ;
    \item Présentation de la méthodologie de travail ;
    \item Explication des objectifs du travail en pointant les idées principales qui seront développées chapitre par chapitre (description de la structure du TFE) ;
\end{itemize} 

Votre TFE fera certainement apparaitre des références. Voici des exemples de références \cite{carter2010back} et \cite{github}. Pour plus d’informations sur comment faire des citations dans des documents LaTeX, vous pouvez vous référer aux slides disponibles sur eCampus (section TFE).

\chapter{État de l’art}

Dresser un État de l'art consiste à : 
\begin{itemize}
    \item Rechercher et présenter toutes les informations existantes concernant le domaine d’étude, la problématique, et/ou les solutions existantes ;
    \item Analyser et comparer les différentes solutions/technologies existantes ;
    \item  En faire une synthèse qui permettra de justifier les choix des différents concepts/ outils/ technologies utilisés dans la solution proposée ;
\end{itemize} 

Attention : l'État de l'art ne consiste pas à remplir le TFE du maximum de matière possible et existante. \textbf{Chaque information ajoutée à l'État de l'art se doit d'être  \underline{pertinente} à la compréhension de la problématique et de la solution développée.} Il est donc important de se poser les questions suivantes : 

\begin{itemize}
    \item \textit{Cette information est-elle en lien avec ma problématique ou bien s'en éloigne-t-elle ?} ;
    \item \textit{Cette information est-elle utile à la compréhension de ma solution ou des choix établis lors de son élaboration ?} ;
\end{itemize} 

En général, l'État de l'art devrait se construire en parallèle de votre implémentation. Il est très déconseillé d'essayer de terminer l'état de l'art avant de commencer à implémenter ou designer la solution étant donné que l'État de l'art est en lien direct avec celle-ci. \footnote{Ceci est une source fréquente d'erreur chez les étudiants.}   

\chapter{Solution(s) proposée(s)}

Cette partie est le cœur du travail de l’étudiant. Elle consiste en la présentation des parties du cycle de développement, un regard critique sur celles-ci et une mise en valeur du travail technique réalisé par l’étudiant. 

\section{Analyse de la solution}

L'analyse dépends du contexte et du travail réalisé. Dans certains cas, il pourra s'agir de diagrammes type UML, Merise, use cases ou encore une présentation des designs patterns utilisés. L'analyse peut autant servir à décrire la solution implémentée qu'à en expliquer sa conception ou à en faire la critique.

\section{Implémentation}

Pour l’implémentation (codage, test et déploiement de l’application, intégration), l’étudiant mettra en valeur le travail réalisé en présentant :

\begin{itemize}
    \item Des métriques pertinentes ;
    \item Des captures d’écran de l’application et/ou des lignes de code spécifiques et
expliquées ;
    \item Des explications concernant toute autre étape/phase qu’il désire mettre en 
évidence, la difficulté rencontrée, le temps mis pour la réalisation, l’originalité
d’un concept/outil utilisé etc... ;
\end{itemize}

\chapter{Limitations et pistes d'amélioration}

Présentez ici les limites de votre solution. Votre solution a très certainement des limitations. Vous avez surement considéré des hypothèses. Il est très peu probable que votre solution soit exempte de défauts. Soyez critique avec votre solution. C'est également l'occasion de dire ce que vous auriez pu améliorer ou faire en plus si vous aviez eu davantage de temps pour développer l'outil.

\chapter{Conclusion}

Elle est généralement organisée comme suit :

\begin{itemize}
    \item Résumé du travail réalisé ;
    \item Rappel des résultats principaux ;
    \item Limitations de la solution proposée ;
    \item Perspectives (pistes pour des améliorations) ;
\end{itemize}

Le texte de la conclusion doit rester neutre, mais doit mettre en avant l’apport de l’étudiant par rapport au sujet.

\chapter{Quelques conseils utiles}

\section*{Complexité technique}

Le TFE a entre autres pour but de montrer la capacité de l’étudiant à mettre en œuvre son savoir et son savoir-faire acquis au cours de sa formation. Votre solution doit donc contenir une complexité technique certaine et un apport personnel suffisant.\\ 

Il peut être difficile de savoir si votre solution est satisfaisante. Le feedback reçu par le promoteur devrait permettre de vous orienter. 

\section*{Introduction et problématique}

L'introduction doit présenter la problématique ainsi que la solution qui va être implémentée. Cette explication de la problématique et de la solution ne doit pas être superficielle. Par exemple, pour la présentation de la solution, dire "\textit{On a créé un outil qui sera présenté dans la partie solution}" n'est pas suffisant. On doit savoir dès l'introduction ce que vous avez développé (même si ce n'est pas dans le détail).\\

Votre problématique doit également être clairement identifiée. Il vous est conseillé de présenter l'introduction sous forme d'une question. Par exemple, un mauvais exemple de présentation de la problématique serait "\textit{Dans ce TFE on va s'intéresser aux outils de Big Data}". Un meilleur exemple serait, "\textit{Ce TFE tentera de répondre à la problématique suivante : Comment les outils de Big Data tel que Hadoop peuvent aider à la visualisation de données dans le contexte d'un site de vente en ligne}".

\section*{Soyez pertinent}

 Les informations présentées doivent être pertinentes. Un problème important souvent rencontré, est que l'état de l'art n'est qu'une agrégation de contenus qui n'ont pas grand-chose en lien direct avec la solution implémentée. Il faut que le contenu présenté serve à comprendre ce qui va être développé en regard avec votre problématique. C'est aussi pour cela qu'il est déconseillé de faire l'état de l'art si vous n'avez encore rien implémenté ou si vous ne savez pas à quoi ressemblera votre solution finale. Evitez donc de remplir votre TFE avec du contenu non pertinent. Le contenu proposé doit permettre une meilleure compréhension de la problématique.   

\section*{Restez objectifs}

Vous n'êtes pas là pour vendre un produit. Evitez les formulations du type "\textit{cet outil est merveilleux pour faire ce travail}", "\textit{cet outil est le plus optimal du marché}", "\textit{c'est outil est vraiment très simple d'utilisation}". Si vous dites qu'un outil est compliqué à mettre en place, vous devez expliquer pourquoi et tenter d'être objectif - idem si vous dites qu'il est simple d'utilisation. Dites-vous qu'on peut toujours vous demander "\textit{Pourquoi pensez-vous que cette technologie est plus simple ?}" et la réponse ne devrait jamais être "\textit{parce que c'est mon avis}" ou "\textit{parce que c'est la société qui le vend qui le dit}". 

\section*{Feedback du promoteur}

Le feedback du promoteur est crucial pour la réalisation du TFE. C'est le premier indicateur pour s'assurer de la bonne réalisation de celui-ci. Il est donc conseillé de présenter régulièrement votre travail, comme préconisé dans le réglement TFE (voir calendrier).

\newpage

\section*{Insertion d'éléments dans LaTeX}

Voici quelques exemples d'insertion d'images et de tableaux dans un document LaTeX. Les images et les tableaux \textbf{doivent toujours contenir une légende}. 

\subsection*{Insertion d'images}

\begin{figure}[h]
    \centering
    \includegraphics[width=0.30\textwidth]{Pictures/GitHub.png}
    \caption{Le logo de GitHub, un service d'hébergement et de gestion de développement \cite{github}.}
    \label{fig:github_logo}
\end{figure}

\subsection*{Insertion de tableaux}

\begin{table}[h!]
\centering
\begin{tabular}{|l|c|c|}
\hline
\rowcolor[HTML]{CBCEFB} 
\textbf{Set} & \multicolumn{1}{l|}{\cellcolor[HTML]{CBCEFB}\textbf{Frequency}} & \multicolumn{1}{l|}{\cellcolor[HTML]{CBCEFB}\textbf{Support}} \\ \hline
A            & 100                                                             & 1.0                                                           \\ \hline
AB           & 50                                                              & 0.5                                                           \\ \hline
ABC          & 25                                                              & 0.25                                                          \\ \hline
\end{tabular}
\caption{Un exemple de tableau généré avec \url{https://www.tablesgenerator.com/}.}
\label{table:1}
\end{table}

%Bibliography
\nocite{*}

\printbibliography

\chapter{Annexes}

\chapter*{Titre de travail de fin d’étude}
\thispagestyle{empty} % allows no page number

Inclure ici un résumé du TFE (minimum 200 mots) sur la quatrième page de couverture
(dernière page au dos du TFE), ainsi que cinq mots clés descriptifs.\\

\Large Mots clés : \normalsize
\begin{itemize}
    \item keyword 0
    \item keyword 1
    \item keyword 2
    \item keyword 3
    \item keyword 4
\end{itemize}

\end{document}
